% \chapter*{Úvod}
% \addcontentsline{toc}{chapter}{Úvod}

% Následuje několik ukázkových kapitol, které doporučují, jak by se
% měla závěrečná práce sázet. Primárně popisují použití \TeX{}ové
% šablony, ale obecné rady poslouží dobře i~uživatelům jiných
% systémů.

\chapter{Úvod}
\addcontentsline{toc}{chapter}{Úvod}

\section*{Motivace}
\addcontentsline{toc}{section}{Motivace}

Deskové hry představují významnou oblast výzkumu umělé inteligence již po několik desetiletí. 
Díky deterministickému prostředí a jednoznačně určeným pravidlům a podmínkám vítězství 
poskytují vhodné prostředí pro návrh, implementaci a porovnávání algoritmů určených 
pro rozhodování. Výzkum v této oblasti vedl ke vzniku různých metod pro prohledávání herních stromů 
a v posledních letech i propojení s technikami ze strojového učení.

Mezi tradiční přístupy pro řešení stavového prostoru deskových her patří algoritmy založené na 
prohledávání herního stromu. Mezi příklady patří algoritmus minimax nebo alfa-beta ořezávání.
Tyto metody jsou stálicí pro hry s relativně omezeným stavovým prostorem, jejich efektivita a
použitelnost ale rychle klesají s rostoucí složitostí hry. V posledních letech se právě proto
do popředí dostaly algoritmy založené na metodě Monte Carlo Tree Search (MCTS), která umožňuje efektivně prohledat
i velký stavový prostor a v kombinaci s neuronovými sítěmi dosahují skvělých výsledků i ve velmi komplexních hrách.

Přestože existuje již nespočet úspěšných algoritmů pro klasické deskové hry, jejich vývoj je velmi úzce spjat 
s konkrétními pravidly dané hry. U nových či méně známých her proto povětšinou neexistuje žádné herní prostředí
ani implementace referenčních botů, díky kterým by bylo možné experimentovat s dalšími algoritmy a s referencí
je pak porovnávat. Vytvoření takového prostředí tedy představuje důležitý předpoklad pro další výzkum a vývoj 
umělé inteligence. 

Tato práce se zaměřuje na návrh a implementaci herního prostředí pro abstraktní deskovou hru inspirovanou hrou s názvem Logix.
Součástí je taktéž vývoj několika herních botů využívajících různé algoritmy pro rozhodování, od jednoduchých 
heuristik až po algoritmus MCTS. Významná část práce je věnována návrhu a implementaci pipeline
pro trénování ohodnocovací funkce založené na neuronové síti. Výsledné dílo má tedy sloužit jako experimentální
platforma pro testování různých dalších přístupů k řešení této hry.


\section*{Cíle práce}
\addcontentsline{toc}{section}{Cíle práce}

Cílem této bakalářské práce je navrhnout a implementovat herní prostředí pro abstraktní deskovou hru inspirovanou hrou Logix,
kde bude možná experimentace s dalšími algoritmy a samotná hra proti implementovaným botům. Součástí je návrh pravidel hry, 
reprezentace herního stavu a implementace prostředí pro simulaci kompletních her.

Další část práce je věnována vývoji několika herních botů využívajících náhodný výběr, heuristiky nebo stromové algoritmy
s hlavním důrazem na MCTS. Ten pak bude sloužit jako hlavní metoda pro výběr tahů.

Významnou součástí bude také implementace trénovací pipeline pro vytvoření ohodnocovací funkce založené na neuronové síti.
Cílem je připravit prostředí pro generování trénovacích dat, architekturu modelu a jeho následné využití pro algoritmus MCTS.

Funkčnost řešení pak bude ověřena pomocí experimentů, které porovnají implementované boty mezi sebou, případně proti
lidským hráčům. Experimenty pak jasně dají najevo, jak velký vliv má natrénovaná ohodnocovací funkce na rozhodování
během hry.

\section*{Struktura práce}
\addcontentsline{toc}{section}{Struktura práce}
